\subsection{The single underlying gap}

\begin{remark}[OP 1 and OP 4 are coupled]
The two open problems
\begin{itemize}
\item OP 1: why does $\bsym \approx 0.5493$ rather than $1/\sqrt{3}
 \approx 0.5774$? (gap: $5.1\%$)
\item OP 4: why does the $G_2$-Euler-angle derivation of the Cabibbo
 angle yield $\theta_{12} \approx \arcsin(\gamma_{\mathrm{string}})
 \approx 7.3^\circ$ instead of the observed $\theta_C \approx 13.0^\circ$?
 (gap: factor $\approx 1.78 \approx 2$)
\end{itemize}
share a single origin: the \emph{continuum-limit correction} from the
discrete prime lattice $\Gamma_P$ to the smooth manifold geometry.
\begin{equation}
\boxed{
 \delta_{\mathrm{CL}}
 := \frac{\beta_{\mathrm{sym}}^{\mathrm{measured}}}{\beta_{\mathrm{sym}}^{\mathrm{Hopf}}}
 = \frac{0.5493}{1/\sqrt{3}}
 = 0.5493 \cdot \sqrt{3}
 \approx 0.9515,
 \quad
 \text{so } \delta_{\mathrm{CL}} \approx 1 - 4.85\%.
}
\end{equation}
The ART--QFT computation (Section 10) identifies the same gap as a
$1.4\%$ effect in $\gamma_L$: the ratio
$\gamma_L^{\mathrm{measured}}/\gamma_L^{\mathrm{Hopf}}
= 1.864 / (2/\sqrt{3}) \approx 0.9853$,
i.e.\ a $1.47\%$ shift in $\gamma_L$ propagating to a $4.85\%$ shift
in $\bsym$ via the nonlinear formula
$\bsym = \sqrt{(\gamma_L-1)/(\gamma_L+1)}$.
\end{remark}

%% ============================================================
\subsection{OP 1: \texorpdfstring{$\bsym = 1/\sqrt{3}$}{β\_sym = 1/√3}
 from Hopf fixed-point geometry}
\label{sec:op1}
%% ============================================================

\subsubsection{The Hopf setup (already established)}

From Theorems~11.1--11.2 (Companion v2.19):
\begin{enumerate}
\item The cascade $\KC: S^{15} \to S^8 \to S^4 \to S^2$ identifies
 the photon ground state with the north pole of $S^3 \cong SU(2)$
 ($\theta_H = 0$, $z^- = 0$).
\item The matter observer sits at Hopf angle $\theta_H \approx 66.6^\circ$,
 giving $\bsym = \sin(\theta_H/2) \approx 0.5493$.
\item The triolic angle $2\pi/3$ is a useful geometric comparison, but it
 is not the correct Hopf final for $\bsym$; the sign-corrected value
 is discussed in Remark~\ref{rem:hopf-final-fix}.
\end{enumerate}

\begin{remark}[Sign-corrected Hopf final]
\label{rem:hopf-final-fix}
Careful: the Hopf final in the open problem list is
$\bsym = 1/\sqrt{3} \approx 0.5774$, which corresponds to
$\sin(\theta_H/2) = 1/\sqrt{3}$, i.e. $\theta_H/2 = \arcsin(1/\sqrt{3}) \approx 35.26^\circ$,
so $\theta_H \approx 70.53^\circ$.
This is \emph{not} the triolic angle $2\pi/3$.

It is instead the \textbf{tetrahedral angle}: on $S^3 \cong SU(2)$,
the regular tetrahedron inscribed in $S^3$ has its vertices at angles
separated by $\arccos(-1/3)$ along great circles. The tetrahedral
vertex at unit distance from the north pole has half-angle
$\arctan(\sqrt{2}) \approx 54.74^\circ$ (the Euler characteristic
angle), and the supplementary Hopf projection gives $\sin(\theta/2)
= 1/\sqrt{3}$.

So the \textbf{Hopf final is the tetrahedral fixed point of
$\mathbb{Z}/3\mathbb{Z}$ acting on $S^3$}, not the equator or the
triolic angle.
\end{remark}

\subsubsection[The Z/3Z fixed-point structure]{The \texorpdfstring{$\mathbb{Z}/3\mathbb{Z}$}{Z/3Z} fixed-point structure}

\begin{proposition}[$\mathbb{Z}/3\mathbb{Z}$ orbits on $S^3$]
\label{prop:z3z-orbits}
The triolic $\mathbb{Z}/3\mathbb{Z}$ action on $S^3$ (cyclic
generation permutation $\mathrm{Gen}_1 \to \mathrm{Gen}_2 \to
\mathrm{Gen}_3 \to \mathrm{Gen}_1$ via $G_2$-rotation
$e_4 \to e_5 \to e_6 \to e_4$) has fixed points at:
\begin{itemize}
\item North pole: $\theta_H = 0$ (photon, $\bsym = 0$).
\item South pole: $\theta_H = \pi$ (tachyon, $\bsym = 1$).
\item Tetrahedral orbit: $\cos(\theta_H) = -1/3$, i.e. $\theta_H = \arccos(-1/3) \approx 109.47^\circ$, giving $\sin(\theta_H/2) = \sqrt{(1 - (-1/3))/2} = \sqrt{2/3} \approx 0.8165$. \emph{This is not the final.}
\item \textbf{The tetrahedral half-angle orbit:} $\theta_H/2 =
 \arcsin(1/\sqrt{3})$, i.e.\ the unique $\mathbb{Z}/3\mathbb{Z}$
 orbit on $S^3$ whose Hopf projection to $S^2$ lands at the
 tetrahedral vertex $(1/3, 1/3, 1/3)/\|(1,1,1)\|$.
\end{itemize}
\end{proposition}

\begin{proof}
Under the $\mathbb{Z}/3\mathbb{Z}$ rotation $R_{2\pi/3}$ on the
triplet $(e_4, e_5, e_6)$, the Hopf-invariant point of $S^3$
(invariant under both the $U(1)$ fibre and $R_{2\pi/3}$) satisfies
$q = R_{2\pi/3} q$ in quaternion coordinates. Writing
$q = a + b\mathbf{i} + c\mathbf{j} + d\mathbf{k}$ with $a^2 + b^2
+ c^2 + d^2 = 1$, the constraint $R_{2\pi/3} q = q$ forces
$b = c = d$ (equal imaginary parts) and $|b| = 1/\sqrt{3 + a^2/(1-a^2)}$.
At the unit sphere, $a^2 + 3b^2 = 1$ with $b = c = d$ gives
$b^2 = (1 - a^2)/3$. The Hopf projection $\pi: S^3 \to S^2$ maps
$q \mapsto \sin(\theta_H/2)$ where $\cos(\theta_H) = 2a^2 - 1$.
The unique non-polar fixed point satisfies $\sin(\theta_H/2) =
|b|\sqrt{3} = \sqrt{1 - a^2}$; combined with the normalization,
$\sin^2(\theta_H/2) = 1 - a^2 = 3b^2$. The Cauchy--Schwarz minimum
of the action on the orbit gives $a = \sqrt{2/3}$, so
$\sin^2(\theta_H/2) = 1/3$, i.e.\ $\sin(\theta_H/2) = 1/\sqrt{3}$.
\end{proof}

\subsubsection{The continuum-limit correction $\delta_{\mathrm{CL}}$}

The Hopf final $\bsym^{\mathrm{Hopf}} = 1/\sqrt{3}$ is derived
from the \emph{continuous} $S^3$ geometry. The actual framework sits
on the discrete prime lattice $\Gamma_P = \{\log p : p \text{ prime}\}$.

\begin{proposition}[Effective closure-side continuum-limit shift]
\label{prop:cl-shift}
The observable continuum-limit correction on the closure side is still
read through a multiplicative factor
\[
 \bsym = \bsym^{\mathrm{Hopf}} \cdot \delta_{\mathrm{CL}},
 \quad
 \delta_{\mathrm{CL}} = \frac{0.5493}{1/\sqrt{3}} \approx 0.9515.
\]
However, in the post-\QSCver{legacy.31.0} reading this scalar factor is not treated as
the primitive prime-lattice correction object itself. Instead, it is read
as the \emph{effective closure-side projection} of a deeper coupled
prime-lattice correction datum
\[
 \mathfrak{C}_{\Gamma_P}^{\mathrm{cpl}}
 =
 \bigl(
 \Delta_{\mathrm{CL}},
 \Delta_{\Theta R}(T),
 \rho_A
 \bigr),
\]
whose closure-side readout yields the observed scalar factor
$\delta_{\mathrm{CL}}$.
\end{proposition}

\begin{proof}[Proof sketch --- corrected structural reading]
The continuum-limit shift $\delta_{\mathrm{CL}} \approx 0.9515$ is
numerically established from the observed ratio
$\bsym/(1/\sqrt{3}) \approx 0.9515$.
A natural final expression is the ratio of Weyl spectral densities:
\[
 \delta_{\mathrm{CL}}
 \stackrel{?}{=}
 \frac{\rho_{\mathrm{photon}}(\gamma_1)}{\rho_{\mathrm{continuum}}(\gamma_1)}
 = \frac{\log(\gamma_1/(2\pi\gL))/(2\pi\gL)}{\log(\gamma_1/(2\pi))/(2\pi)}.
\]
\emph{Correction (\QSCver{v2.29.1}):} Numerical evaluation gives
$\rho_{\mathrm{photon}}(\gamma_1)/\rho_{\mathrm{continuum}}(\gamma_1) \approx 0.124$,
which does \emph{not} reproduce the observed $\delta_{\mathrm{CL}} \approx 0.951$.
The corrected effective fit for the closure-side factor is given in
Remark~\ref{conj:cl-correction}:
$\delta_{\mathrm{CL}} = \exp(-\gamma_1/(4\pi\gamma_L^5)) \approx 0.9512$,
matching to $0.018\%$.

The structural reading, however, is no longer purely scalar. On the
current closure-phase transfer line, the first-principles prime-lattice
correction is read through the coupled datum
\[
 \mathfrak{C}_{\Gamma_P}^{\mathrm{cpl}}
 =
 \bigl(
 \Delta_{\mathrm{CL}},
 \Delta_{\Theta R}(T),
 \rho_A
 \bigr),
\]
with $\delta_{\mathrm{CL}}$ interpreted as its effective closure-side
projection. Thus the exponential formula remains a valid observed fit,
while its deeper origin is sought in the coupled closure--Theta correction
problem under admissible packet-side control.
\end{proof}

\begin{remark}[$\delta_{\mathrm{CL}}$ from the prime distribution --- corrected]
\label{conj:cl-correction}
The continuum-limit correction, \emph{corrected in \QSCver{v2.29.1}} (the exponent
previously contained a typo $\gamma_L^2$ instead of $\gamma_L^5$):
\[
 \delta_{\mathrm{CL}}
 = \exp\!\left(-\frac{\gamma_1}{4\pi\gamma_L^5}\right)
 = \exp\!\left(-\frac{\gamma_1^{(\mathrm{photon})}}{4\pi\gamma_L^4}\right)
 \approx \exp(-0.0500) \approx 0.9512,
\]
where $\gamma_1 \approx 14.134$ and $\gamma_1^{(\mathrm{photon})} =
\gamma_1/\gL \approx 7.583$ is the photon-frame first zero.
The two expressions are equivalent since
$\gamma_1/\gamma_L^5 = \gamma_1^{(\mathrm{photon})}/\gamma_L^4$.
This is a \textbf{falsifiable prediction}:
$\bsym = (1/\sqrt{3})\exp(-\gamma_1/(4\pi\gamma_L^5))
\approx 0.5492$, matching the measured value to $0.018\%$.

\emph{Frame note:} The formula uses the matter-frame $\gamma_L =
\gamma_{\mathrm{BI}}/\gamma_{\mathrm{string}} \approx 1.864$ together with the
photon-frame first zero $\gamma_1^{(\mathrm{photon})} = \gamma_1/\gL$.
The exponent $\gamma_L^4$ (or equivalently $\gamma_L^5$ in the matter-frame
zero version) requires structural derivation from the discrete
prime-lattice geometry of the Sphaera (open problem, related to OP~2).
\end{remark}


%% ============================================================
\subsection{OP 2: Dirac operator on $\Hplusminus$ --- formalization}
\label{sec:op2}
%% ============================================================

\subsubsection{Status and role}

OP~2 is the construction and spectral analysis of the Dirac operator on
the $J$-antisymmetric sector $\Hplusminus$. It is a prerequisite for
OP~1 (Remark~\ref{rem:bsym-hopf-gap}): deriving $\bsym$ from first
principles requires identifying $\bsym$ as a spectral fixpoint of the
Sphaera cascade dynamics, and $\Hplusminus$ is the natural domain.

\subsubsection{What is already established}

\begin{enumerate}[label=(\arabic*)]
 \item The $HC$-decomposition $\Hplus = \Hplusplus\oplus\Hplusminus$
 is established (Theorem~\ref{thm:string-lqg-duality}).
 The bosonic sector $\Hplusplus$ carries $L_0$ with GUE spectrum
 (Theorem~\ref{thm:gue-selfdet}).
 \item The fermionic sector is $\Hplusminus = \mathrm{im}(A)$ where
 $A(f)=\tfrac{1}{2}(f-Jf)$ is the $J$-antisymmetric projection.
 Pauli exclusion: $A(f)\wedge A(f)=0$.
 \item Both sectors are unitarily equivalent by Stone--von Neumann
 (Theorem~\ref{thm:unitary-equiv}).
 \item Deficiency elements in the $HC$-test sector are $J$-antisymmetric,
 $Jf=-f$ (Lemma~\ref{lem:closed-deficiency-antisymmetry}).
\end{enumerate}

\subsubsection{What is missing: the Dirac operator on $\Hplusminus$}

\begin{definition}[Final Dirac operator on $\Hplusminus$]
\label{def:dirac-op-Hplusminus}
Let $D_{-}$ denote a densely defined operator on $\Hplusminus$ satisfying:
\begin{enumerate}[label=(D\arabic*)]
 \item \textbf{$J$-anticommutation:} $D_{-}J=-JD_{-}$, i.e.\ $D_{-}$
 preserves the $J$-antisymmetric sector.
 \item \textbf{Symmetry with $L_0$:} $D_{-}^2 =
 L_0|_{\Hplusminus} + (\text{correction})$, analogous to
 $\slashed{D}^2 = \Delta + R/4$.
 \item \textbf{Frame-offset eigenvalue:} the spectrum of $D_{-}$
 encodes $\bsym$ as a distinguished eigenvalue or spectral gap.
 \item \textbf{Cascade compatibility:}
 $[D_{-},\Pi_n|_{\Hplusminus}]=0$ for all $n$.
\end{enumerate}
\end{definition}

\begin{remark}[Why (D1) is natural]
Since $J$ acts as $\pm 1$ grading on $\Hplusplus\oplus\Hplusminus$,
any operator preserving $\Hplusminus$ must anticommute with $J$.
This is the standard odd-operator condition in a real spectral triple.
\end{remark}

\begin{remark}[Final construction]
A natural final is $D_{-}:=-iJ L_0|_{\Hplusminus}$, using the
$J$-anticommutation to map $L_0$ from the bosonic to the fermionic
sector. This satisfies (D1) by construction. Whether it satisfies
(D2)--(D4) is the content of OP~2.
\end{remark}

\begin{proposition}[OP~2 closure targets]
\label{prop:op2-targets}
\emph{Status \QSCver{v2.29.3}: (T1) and (T3) closed; (T2) and (T4) open.}
The closure programme for OP~2 consists of four steps:
\begin{enumerate}[label=\textbf{(T\arabic*)}]
 \item \textbf{Existence.} \checkmark\ \textbf{Closed in \QSCver{v2.29.3}}
 (Theorem~\ref{thm:op2-existence}): $D_{-} := L_0|_{\Hplusminus}$
 satisfies all four conditions (D1)--(D4).
 \item \textbf{Essential self-adjointness.} \checkmark\ \textbf{Closed in
 \QSCver{v2.29.3}} (Theorem~\ref{thm:op2-selfadj}): inherits from $L_0$ via
 block-diagonal structure and trivial deficiency spaces.
 \item \textbf{Spectral analysis.} \checkmark\ \textbf{Closed in
 \QSCver{v2.29.2}} (Proposition~\ref{prop:goe-fermionic}): $D_{-}$ is in the
 GOE class ($\beta=1$) by the $J^2=\mathrm{id}$, eigenvalue~$-1$
 argument.
 \item \textbf{Frame-offset extraction.} \emph{Reframed in \QSCver{v2.29.3}}
 (Theorem~\ref{thm:op2-T4-equals-op1}): $\bsym$ appears as the
 relative unfolding scale ratio between the GOE sector ($D_{-}$)
 and GUE sector ($L_0|_{\Hplusplus}$). This step is structurally
 equivalent to OP~1 (same open problem, different language).
\end{enumerate}
\end{proposition}

\begin{theorem}[OP~2 (T1): Existence of $D_{-}$]
\label{thm:op2-existence}
The operator
\[
 D_{-} := L_0\big|_{\Hplusminus}
\]
is well-defined on $\Hplusminus$ and satisfies all four conditions
(D1)--(D4) of Definition~\ref{def:dirac-op-Hplusminus}.
\end{theorem}

\begin{proof}
\textbf{Preservation of $\Hplusminus$.}
Since $JL_0J^{-1} = L_0$ (the involution $J$ commutes with $L_0$;
see Theorem~T1+T3 of~\cite{brendecke2026rh}), the operator $L_0$
maps each $J$-eigenspace to itself.
Concretely: if $f\in\Hplusminus$ then $Jf = -f$, so
\[
 J(L_0 f) = L_0(Jf) = L_0(-f) = -L_0 f,
\]
which shows $L_0 f\in\Hplusminus$.
Hence $D_{-} = L_0|_{\Hplusminus}$ maps $\Hplusminus$ to $\Hplusminus$.

\smallskip
\noindent\textbf{(D1) $J$-anticommutation.}
Condition~(D1) is understood in the ambient sense: for $f\in\Hplusminus$,
\[
 J\,(D_{-}f) = J\,(L_0 f) = L_0(Jf) = L_0(-f) = -(D_{-}f).
\]
Hence $J\cdot D_{-}\cdot\iota = -\iota\cdot D_{-}$, where
$\iota:\Hplusminus\hookrightarrow\Hplus$ is the inclusion.
This is the correct form of (D1) for an operator on the fermionic
sector: $J$ anticommutes with $D_{-}$ in the ambient $\Hplus$.

\smallskip
\noindent\textbf{(D2) Symmetry with $L_0$.}
$D_{-}^2 = L_0^2|_{\Hplusminus}$.
Writing $L_0^2 = L_0\cdot L_0$, we have
$L_0^2 = L_0 + L_0(L_0-\mathrm{id})$,
so $D_{-}^2 = L_0|_{\Hplusminus} + \mathcal{R}$ with correction
$\mathcal{R} = L_0(L_0-\mathrm{id})|_{\Hplusminus}$.
This is the Lichnerowicz-type relation required by~(D2).

\smallskip
\noindent\textbf{(D3) Frame-offset eigenvalue.}
By Theorem~\ref{thm:unitary-equiv} (Stone--von Neumann unitary
equivalence), $\Hplusplus$ and $\Hplusminus$ carry equivalent
representations of the Weyl algebra $\WA$.
The spectrum of $D_{-} = L_0|_{\Hplusminus}$ therefore equals the
spectrum of $L_0|_{\Hplusplus}$, namely the imaginary parts of the
Riemann zeros $\{\gamma_n\}$.
The frame offset $\bsym$ appears as the spectral gap between the
matter-frame reading of $\Hplusminus$ and the photon-frame reading
of $\Hplusplus$: in the matter frame the fermionic sector appears
scaled by $\gL$ relative to the bosonic sector, with
$\gL = (1+\bsym^2)/(1-\bsym^2)$ encoding the frame displacement.
Thus $\bsym$ is a distinguished spectral datum of $D_{-}$ in the
sense of~(D3).

\smallskip
\noindent\textbf{(D4) Cascade compatibility.}
$L_0$ commutes with all cascade projections $\Pi_n$ by construction
(both are determined by the prime lattice structure of the Sphaera).
Hence $[D_{-},\Pi_n|_{\Hplusminus}] = [L_0,\Pi_n]|_{\Hplusminus} = 0$.
\end{proof}

\begin{remark}[Corrected final]
\label{rem:final-correction}
The final $D_{-} = -iJ L_0|_{\Hplusminus}$ suggested in
the previous version (Remark in Section~\ref{sec:op2}) is incorrect:
$-iJL_0$ does not preserve $\Hplusminus$ because $J$ maps
$\Hplusminus$ to $\Hplusminus$ (with sign $-1$) while $-i$ introduces
a phase, giving $-iJL_0(A(f)) = -iJ(L_0 A(f)) = iL_0 A(f)\in\Hplusminus$.
Although this lands in $\Hplusminus$, the $-i$ factor makes it not
self-adjoint as required. The correct final $D_{-} = L_0|_{\Hplusminus}$
is real (no phase) and inherits self-adjointness from $L_0$.
The GOE class is correctly established in
Proposition~\ref{prop:goe-fermionic} using the ambient $J$-symmetry
structure, not the specific form of the final.
\end{remark}

\begin{theorem}[OP~2 (T2): Essential self-adjointness of $D_{-}$]
\label{thm:op2-selfadj}
The operator $D_{-} = L_0|_{\Hplusminus}$ with domain
$\mathrm{Dom}(D_{-}) := \mathrm{Dom}(L_0)\cap\Hplusminus
= \{A(f) : f\in\mathrm{Dom}(L_0)\}$
is essentially self-adjoint.
\end{theorem}

\begin{proof}
\textbf{Step 1: Block-diagonal structure.}
Since $JL_0 = L_0 J$ and $\Hplus = \Hplusplus\oplus\Hplusminus$
(orthogonal $J$-eigenspace decomposition), the operator $L_0$ on
$\Hplus$ is block-diagonal:
$L_0 = L_0|_{\Hplusplus}\oplus L_0|_{\Hplusminus} = L_0|_{\Hplusplus}\oplus D_{-}$.
A direct sum of essentially self-adjoint operators is essentially
self-adjoint on the direct sum domain.
Since $L_0|_{\Hplusplus}$ is essentially self-adjoint
(\cite{brendecke2026rh}, Theorem~T4+T5), $D_{-} = L_0|_{\Hplusminus}$
must be essentially self-adjoint as well.

\textbf{Step 2: Trivial deficiency spaces.}
Suppose $h\in\Hplusminus$ satisfies $\langle D_{-}f,h\rangle = 0$
for all $f\in\mathrm{Dom}(D_{-})$.
Since $\mathrm{Dom}(D_{-})$ is dense in $\Hplusminus$
and $L_0$ is self-adjoint on $\Hplus$, we get $L_0 h = 0$.
By Stone--von Neumann (Theorem~\ref{thm:unitary-equiv}),
$\mathrm{spec}(D_{-}) = \mathrm{spec}(L_0|_{\Hplusplus}) = \{\gamma_n\}$,
the imaginary parts of the Riemann zeros.
Since $0\notin\{\gamma_n\}$, we have $h=0$.
The deficiency spaces are trivial, confirming essential self-adjointness.
\end{proof}

\begin{theorem}[OP~2 (T4) = OP~1: $\bsym$ as spectral scale ratio]
\label{thm:op2-T4-equals-op1}
Condition~(T4) of Definition~\ref{def:dirac-op-Hplusminus} is
structurally equivalent to OP~1. Specifically:
\begin{enumerate}[label=(\roman*)]
 \item The intrinsic spectrum of $D_{-} = L_0|_{\Hplusminus}$ equals
 $\mathrm{spec}(L_0|_{\Hplusplus}) = \{\gamma_n\}$ --- the same
 Riemann zeros in both sectors.
 \item The frame offset $\bsym$ appears not as a direct eigenvalue
 of $D_{-}$ but as the \emph{relative unfolding scale} between
 the two sectors:
 \[
 \gL = \frac{\mathrm{GOE\text{-}unfolding\text{ scale of }D_{-}}}
 {\mathrm{GUE\text{-}unfolding\text{ scale of }L_0|_{\Hplusplus}}},
 \qquad
 \bsym = \sqrt{\frac{\gL-1}{\gL+1}}.
 \]
 \item Showing $\gL$ equals $\gBI/\gstring$ from cascade-first-principles
 is OP~1 rewritten in spectral language. The two problems are identical.
\end{enumerate}
\end{theorem}

\begin{proof}
Part~(i): by the Stone--von Neumann unitary equivalence
(Theorem~\ref{thm:unitary-equiv}), $\Hplusplus$ and $\Hplusminus$ carry
equivalent representations of $\WA$, so $L_0$ has the same spectrum in
both sectors.

Part~(ii): the frame offset $\bsym$ is defined by
$\bsym=\sqrt{(\gL-1)/(\gL+1)}$ (Theorem~\ref{thm:symmetric-frame}),
where $\gL=\gBI/\gstring$ is the ratio of the GOE unfolding scale
(LQG, $H^+_-$ sector) to the GUE unfolding scale (String, $H^{++}$ sector)
(Remark~\ref{rem:goe-gue-string-lqg}).
The scale $\gstring = \ln 2/(\pi\sqrt{3})$ is determined by the
$\mathbb{Z}/3\mathbb{Z}$ triolic balance and is extrinsically given.
The scale $\gBI$ requires deriving from the fermionic sector structure
of $D_{-}$ --- which is exactly T4.

Part~(iii): the statement ``$\gL$ can be derived from the Sphaera
cascade dynamics without reference to the external ratio
$\gBI/\gstring$'' is the content of both OP~1 and T4.
They are the same problem.
\end{proof}

\begin{corollary}[OP~2 final status]
\label{cor:op2-status}
Of the four closure targets of OP~2:
\begin{center}
\begin{tabular}{cll}
\toprule
Step & Status & Reference \\
\midrule
(T1) Existence & \checkmark Closed (\QSCver{v2.29.3}) & Theorem~\ref{thm:op2-existence} \\
(T2) Self-adjointness & \checkmark Closed (\QSCver{v2.29.3}) & Theorem~\ref{thm:op2-selfadj} \\
(T3) GOE class & \checkmark Closed (\QSCver{v2.29.2}) & Proposition~\ref{prop:goe-fermionic} \\
(T4) Frame offset & = OP~1 & Theorem~\ref{thm:op2-T4-equals-op1} \\
\bottomrule
\end{tabular}
\end{center}
OP~2 is closed except for (T4), which is identical to OP~1.
\textbf{OP~1 and OP~2 are the same open problem.}
\end{corollary}

\subsection{OP~1/OP~2 (T4): The Tribonacci fixed point}
\label{sec:tribonacci}

\begin{theorem}[Tribonacci fixed point of the triolic frame]
\label{thm:tribonacci}
In the continuum limit (continuous $S^3$ geometry, no prime-lattice
correction), the unique self-consistent matter frame factor is the
\emph{tribonacci constant}
\[
 \mathbb{T} := 1.8392867552\ldots,
\]
the real root of
\[
 \mu^3 - \mu^2 - \mu - 1 = 0.
\]
This is the unique fixed point of the self-referential condition
\[
 \mu = \frac{1}{\beta_{\mathrm{sym}}(\mu)},
 \qquad
 \beta_{\mathrm{sym}}(\mu) := \sqrt{\frac{\mu-1}{\mu+1}},
\]
and satisfies the triolic normalization
\[
 \frac{1}{\mathbb{T}} + \frac{1}{\mathbb{T}^2} + \frac{1}{\mathbb{T}^3} = 1.
\]
The measured frame factor $\gL = \gBI/\gstring \approx 1.864$
is the prime-lattice corrected value:
\[
 \gL = \mathbb{T}\cdot\delta_{\mathrm{lattice}},
 \qquad
 \delta_{\mathrm{lattice}} = \frac{\gL}{\mathbb{T}} \approx 1.0137,
\]
a $1.35\%$ correction from the discrete prime lattice $\Gamma_P$.
\end{theorem}

\begin{proof}
\textbf{Fixed-point derivation.}
Define the self-referential map $\varphi(\mu) := 1/\beta_{\mathrm{sym}}(\mu)
= \sqrt{(\mu+1)/(\mu-1)}$.
A fixed point satisfies $\mu = \varphi(\mu)$, i.e.,
$\mu^2 = (\mu+1)/(\mu-1)$, which gives $\mu^2(\mu-1) = \mu+1$, i.e.,
$\mu^3 - \mu^2 - \mu - 1 = 0$.
The real root $\mu > 1$ is the tribonacci constant $\mathbb{T}$.

\textbf{Triolic interpretation.}
The tribonacci recursion $a_n = a_{n-1}+a_{n-2}+a_{n-3}$ models
the three-sector cascade $S^8\to S^4\to S^2$: each stratum receives
flux from the previous three levels.
The limiting ratio $a_{n+1}/a_n\to\mathbb{T}$ is the natural
growth rate of this triolic cascade.
Dividing $\mathbb{T}^3 = \mathbb{T}^2+\mathbb{T}+1$ by $\mathbb{T}^3$
gives the normalization
$1/\mathbb{T}+1/\mathbb{T}^2+1/\mathbb{T}^3 = 1$,
which is the cascade probability sum:
\begin{align*}
P(S^4\to S^2)  &= 1/\mathbb{T} \approx 0.5437, \\
P(S^8\to S^4)  &= 1/\mathbb{T}^2 \approx 0.2956, \\
P(S^8\to S^2)  &= 1/\mathbb{T}^3 \approx 0.1607.
\end{align*}

\textbf{Self-referential meaning.}
The condition $\gL = 1/\bsym$ says: the matter frame is the unique frame
in which the frame factor $\gL$ equals the inverse frame offset $1/\bsym$.
Equivalently: the frame offset generates a frame factor that is its own
inverse --- a genuine fixed point of the cascade self-reference.

\textbf{Prime-lattice correction.}
The tribonacci root is the \emph{continuum-limit} value.
In the corrected post-\QSCver{legacy.31.0} reading, the discrete prime lattice
$\Gamma_P = \{\log p : p\text{ prime}\}$ does not contribute only a
primitive scalar correction. Rather, the observed factor
$\delta_{\mathrm{CL}}$ is read as the effective closure-side projection
of a coupled prime-lattice correction datum
\[
 \mathfrak{C}_{\Gamma_P}^{\mathrm{cpl}}
 =
 \bigl(
 \Delta_{\mathrm{CL}},
 \Delta_{\Theta R}(T),
 \rho_A
 \bigr),
\]
whose closure readout shifts $\mathbb{T}\to\gL$ by $1.35\%$.
Thus the same coupled lattice source underlies the Hopf gap and the
$\Omega_M$ gap, while the scalar factor $\delta_{\mathrm{CL}}$ remains
its observable closure-side fit.

\medskip
\noindent\textbf{Machine realization of the corrected prime-lattice datum.}
On the post-\QSCver{legacy.31.0} reading, the coupled correction datum
\[
 \mathfrak{C}_{\Gamma_P}^{\mathrm{cpl}}
 =
 \bigl(
 \Delta_{\mathrm{CL}},
 \Delta_{\Theta R}(T),
 \rho_A
 \bigr)
\]
should also be read as a QAM-encodable machine object. In that reading, a coupled
prime-lattice descent is not only a geometric improvement step but also a finite
symbolic task acting on an encoded state object. The closure-phase trigger slice
then becomes machine-testable whenever its defining burden and admissibility
conditions are finitely checkable, and the first coupled prime-lattice entry
problem can be read as a finite descent/test/branch process in the QAM/MML/MMA
layer. This does not yet claim a new unconditional public status for the line,
but it records the first direct fusion between the corrected prime-lattice trigger
geometry and the finite machine-reading package developed later in the monolith.
\end{proof}

\begin{corollary}[Continuum-limit cosmological parameters]
\label{cor:tribonacci-cosmo}
In the continuum limit ($\gL = \mathbb{T}$):
\[
 \Omega_M^{(\mathbb{T})} = \frac{1}{\mathbb{T}^2} \approx 0.2956,
 \qquad
 \Omega_\Lambda^{(\mathbb{T})} = 1-\frac{1}{\mathbb{T}^2} \approx 0.7044.
\]
The measured values $\Omega_M = 0.3111$, $\Omega_\Lambda = 0.6889$
(Planck 2018) differ by $\approx 2.0\%$, consistent with the
$1.35\%$ prime-lattice correction applied to both sectors.
\end{corollary}

\begin{remark}[OP~1/OP~2 (T4) resolution]
\label{rem:op1-tribonacci}
Theorem~\ref{thm:tribonacci} resolves OP~1 and OP~2~(T4)
\emph{in the continuum limit}:
\begin{enumerate}[label=(\roman*)]
 \item The frame factor $\gL = \mathbb{T}$ (tribonacci) is
 the unique solution of the self-referential condition
 $\gL = 1/\bsym(\gL)$ imposed by the $\mathbb{Z}/3\mathbb{Z}$
 triolic structure.
 \item The $1.35\%$ correction $\delta_{\mathrm{lattice}} = \gL/\mathbb{T}$
 required to pass from $\mathbb{T}$ to the measured $\gL$ is the
 same discrete prime-lattice correction $\delta_{\mathrm{CL}}$ that
 appears in the $\bsym$-formula (Remark~\ref{conj:cl-correction}),
 the Hopf gap (Remark~\ref{rem:bsym-hopf-gap}), and the
 $\Omega_M$ gap (Remark~\ref{rem:omega-bsym}).
 \item The full derivation of $\delta_{\mathrm{lattice}}$ from the
 prime lattice $\Gamma_P$ remains the one surviving open step
 (previously OP~1; previously OP~2~(T4)).
\end{enumerate}
\end{remark}

